\documentclass{szb-solution}

\area{Kalkulus}
\title{Integrálszámítás I}
\subject{Matematika G1}
\subjectCode{BMETE94BG01}
\date{Utoljára frissítve: \today}
\docno{10}

\begin{document}
\maketitle

\subsection{Az első függvény teljes vizsgálata}

Az értelmezési tartomány $\Domain_f=\mathbb R\setminus\{0\}$. A függvénynek
nincs zérushelye, nem páros, nem páratlan és nem periodikus. A fontos
határértékek:
\begin{align*}
  \lim_{x\to-\infty}f(x)&=-\infty,
  &\lim_{x\to+\infty}f(x)&=+\infty,\\
  \lim_{x\to0^-}f(x)&=-\infty,
  &\lim_{x\to0^+}f(x)&=0
  \text.
\end{align*}

Az első derivált
$$
  f'(x)=e^{-1/x}\left(1+\frac1x\right)
       =e^{-1/x}\frac{x+1}{x}
  \text.
$$
Mivel az exponenciális tényező pozitív, $f$ szigorúan növekvő a
$(-\infty,-1)$ és $(0,\infty)$ intervallumokon, szigorúan csökkenő a
$(-1,0)$ intervallumon. Az $x=-1$ pontban lokális maximuma van:
$$
  \boxed{f(-1)=-e}
  \text.
$$

A második derivált különösen egyszerű:
$$
  f''(x)=\frac{e^{-1/x}}{x^3}
  \text.
$$
Ezért a függvény konkáv a $(-\infty,0)$, konvex a $(0,\infty)$
intervallumon. A görbületváltás helye, $x=0$, nem tartozik az értelmezési
tartományhoz, tehát nincs inflexiós pont.

Az $x=0$ egyenes bal oldali függőleges aszimptota. A végtelenben ferde
aszimptotát keresünk. Mivel
$$
  \lim_{x\to\pm\infty}\frac{f(x)}x=1
$$
és
$$
  \lim_{x\to\pm\infty}\bigl(f(x)-x\bigr)
  =\lim_{x\to\pm\infty}x\bigl(e^{-1/x}-1\bigr)=-1,
$$
a ferde aszimptota
$$
  \boxed{y=x-1}
  \text.
$$
A monotonitás és a határértékek alapján
$$
  \boxed{\Range_f=(-\infty,-e]\cup(0,\infty)}
  \text.
$$

\subsection{A második függvény teljes vizsgálata}

A függvény minden valós számon értelmezett, $2\pi$ szerint periodikus, de
sem nem páros, sem nem páratlan. Hasznos alakja
$$
  g(x)=(\sin x-1)^2-1
  \text,
$$
amelyből rögtön látszik, hogy $\Range_g=[-1,3]$. A zérushelyek
$$
  \sin x(\sin x-2)=0,
  \qquad
  \boxed{x=k\pi\quad(k\in\mathbb Z)}
  \text.
$$

Az első derivált
$$
  g'(x)=2\cos x(\sin x-1)
  \text.
$$
Egy periódus alapján, minden $k\in\mathbb Z$ esetén $g$ szigorúan
\begin{itemize}
  \item csökken a
        $(-\pi/2+2k\pi,\,\pi/2+2k\pi)$ intervallumon;
  \item nő a
        $(\pi/2+2k\pi,\,3\pi/2+2k\pi)$ intervallumon.
\end{itemize}
A szélsőértékek:
$$
  \boxed{g(\pi/2+2k\pi)=-1\quad\text{(minimum)}}
$$
és
$$
  \boxed{g(3\pi/2+2k\pi)=3\quad\text{(maximum)}}
  \text.
$$

A második derivált
$$
  g''(x)=-4\sin^2x+2\sin x+2
        =-2(2\sin x+1)(\sin x-1)
  \text.
$$
A $\sin x=1$ megoldásokban a második derivált nulla, de nem vált előjelet,
ezért ezek nem inflexiós pontok. Az előjel a $\sin x=-1/2$ helyeken vált,
így az inflexiós helyek
$$
  \boxed{x=\frac{7\pi}{6}+2k\pi
  \quad\text{és}\quad
  x=\frac{11\pi}{6}+2k\pi}
  \qquad(k\in\mathbb Z)
  \text.
$$
A függvény konkáv ott, ahol $\sin x<-1/2$, és konvex ott, ahol
$\sin x>-1/2$; a fenti elszigetelt $\sin x=1$ pontok ezt nem változtatják
meg. Nem állandó periodikus függvényként nincs lineáris aszimptotája.

\subsection{Határozatlan integrálok}

Minden eredményben $C\in\mathbb R$ tetszőleges integrációs állandó.

\begin{enumerate}[a)]
  \item A gyökös tagot hatványalakba írjuk:
        $$
          \sqrt{x\sqrt{x\sqrt{x}}}=x^{7/8}
          \qquad(x\geq0)
          \text.
        $$
        Ezért
        \begin{align*}
          \int\left[x^2(x^2-1)
          +2\sqrt{x\sqrt{x\sqrt{x}}}\right]\dd x
          &=\int\left(x^4-x^2+2x^{7/8}\right)\dd x\\
          &=\boxed{\frac{x^5}{5}-\frac{x^3}{3}
          +\frac{16}{15}x^{15/8}+C}
          \text.
        \end{align*}

  \item Polinomos osztással
        $$
          \frac{x^2-4x+7}{x-2}=x-2+\frac3{x-2}
          \text,
        $$
        így
        $$
          \boxed{\int\frac{x^2-4x+7}{x-2}\dd x
          =\frac{x^2}{2}-2x+3\ln|x-2|+C}
          \text.
        $$

  \item A $\tan^2x=1/\cos^2x-1$ azonosságot használva
        $$
          \boxed{\int\tan^2x\dd x=\tan x-x+C}
          \text.
        $$

  \item Először egyszerűsítünk:
        $$
          \frac{e^{3x}+1}{e^x+1}=e^{2x}-e^x+1
          \text.
        $$
        Továbbá
        $1/\sqrt{2+2x^2}=1/(\sqrt2\sqrt{1+x^2})$. Így
        \begin{align*}
          &\int\left(
          \frac{\ln2}{\sqrt{2+2x^2}}
          +\frac{e^{3x}+1}{e^x+1}\right)\dd x\\
          &\qquad=\boxed{
          \frac{\ln2}{\sqrt2}\operatorname{arsinh}x
          +\frac12e^{2x}-e^x+x+C}
          \text.
        \end{align*}
        Az első tag helyett
        $\frac{\ln2}{\sqrt2}\ln|x+\sqrt{x^2+1}|$ is írható.
\end{enumerate}

\subsection{Integrálás helyettesítéssel}

\begin{enumerate}[a)]
  \item Legyen $u=x^4+4x^2$. Ekkor
        $\dd u=4(x^3+2x)\dd x$, ezért
        $$
          \boxed{\int(x^3+2x)\cos(x^4+4x^2)\dd x
          =\frac14\sin(x^4+4x^2)+C}
          \text.
        $$

  \item Legyen $u=\tan x$, ekkor $\dd u=\dd x/\cos^2x$. Tehát
        $$
          \boxed{\int\frac{\sqrt[3]{\tan x}}{\cos^2x}\dd x
          =\frac34(\tan x)^{4/3}+C}
          \text.
        $$

  \item Az $u=x^2$ helyettesítéssel $\dd u=2x\dd x$, ezért
        $$
          \boxed{\int\frac{x}{x^4+1}\dd x
          =\frac12\arctan(x^2)+C}
          \text.
        $$

  \item Legyen $u=\sin(2x+1)$, így
        $\dd u=2\cos(2x+1)\dd x$. Ebből
        $$
          \boxed{\int\sin^3(2x+1)\cos(2x+1)\dd x
          =\frac18\sin^4(2x+1)+C}
          \text.
        $$

  \item Az $u=4+x^2$ helyettesítésből $\dd u=2x\dd x$, tehát
        $$
          \boxed{\int\frac{x}{4+x^2}\dd x
          =\frac12\ln(4+x^2)+C}
          \text.
        $$

  \item Az $u=\tan x$ helyettesítéssel
        $\dd u=\dd x/\cos^2x$, valamint
        $\sin x\cos x=u\cos^2x$. Ezért
        $$
          \boxed{\int\frac{1}{\sin x\cos x}\dd x
          =\ln|\tan x|+C}
          \text.
        $$

  \item Mivel $1/\tan x=\cos x/\sin x$, az $u=\sin x$ helyettesítés adja:
        $$
          \boxed{\int\frac1{\tan x}\dd x=\ln|\sin x|+C}
          \text.
        $$

  \item Legyen $u=\ln x$, ahol $x>0$. Ekkor $\dd u=\dd x/x$, így
        $$
          \boxed{\int\frac1{x\ln x}\dd x=\ln|\ln x|+C}
          \text.
        $$

  \item A számlálót kettébontjuk:
        \begin{align*}
          \int\frac{2x+1}{x^2+2}\dd x
          &=\int\frac{2x}{x^2+2}\dd x
            +\int\frac1{x^2+2}\dd x\\
          &=\boxed{\ln(x^2+2)
            +\frac1{\sqrt2}\arctan\!\left(\frac{x}{\sqrt2}\right)+C}
          \text.
        \end{align*}
\end{enumerate}

\end{document}
